% !TEX root = ../main.tex
\chapter{Tensor space theory II: over a ring}\label{ch:l11}
\begin{paracol}{2}

\begin{sources}
\begin{tabularx}{\linewidth}{@{}l l L@{}}
\toprule
\textbf{Video} & \textbf{Length} & \textbf{Content}\\
\midrule
\seg{11}{1}{Lecture 11: Tensor space theory II: over a ring} & 1:55:10 & Vector fields as smooth sections of $TM$; the pointwise addition and the multiplication by functions; rings and the ring $C^\infty(M)$; modules over a ring; modules over a division ring have a basis (Zorn's lemma, with the student question on upper bounds),
the hairy ball as an example without a basis; free and projective modules; the theorem on sections of vector bundles; $\mathrm{Hom}$ of projective modules over a commutative ring; the textbook definition of tensor fields.\\
\bottomrule
\end{tabularx}

\smallskip
\textbf{Transcript} \file{materials/transcripts/C11-V01.txt} (YouTube automatic captions).
\end{sources}

\begin{reading}[Where this lecture sits in the other notes]
\emph{Gravity and Light}, Lecture~6 (Fields): vector fields and the $C^\infty(M)$-module structure of the set of fields; \emph{Theoretische Mechanik}, Chapter~5 (tensor fields). Both treat vector and tensor fields on a smooth manifold; the emphasis here on rings, modules, Zorn's lemma and the absence of a global basis
is, as far as the titles say, special to this course. (From the lecture titles and the plan of the course; not checked in detail.)
\end{reading}

% ==================================================================
\section{Vector fields and the algebraic structure they carry}\label{sec:rm-fields}

Until now (\cref{ch:l09,ch:l10}) the tangent space of a manifold was considered only at one point; last time the tangent bundle, the disjoint union of all the tangent spaces, was constructed, and it is a \emph{smooth} bundle. That now allows one to define vector fields and, in general, tensor fields. Over one point the tangent space is a vector
space; but for the \emph{collection} of vector fields the structure will turn out to be only a \textbf{module}, not a vector space. ``Most people who apply mathematics consider rings and modules a somewhat esoteric subject'': today it will be seen that the funny properties of modules, compared with vector spaces, are of direct geometric
relevance and make one understand the subject better \ts{11}{1}{0:04} \ts{11}{1}{1:06} \ts{11}{1}{1:31} \ts{11}{1}{1:51}.

\begin{definition}[New concepts: vector field, $\Gamma(TM)$]\label{def:rm-field}
Let $M$ be a smooth manifold and $\pi\colon TM\to M$ its tangent bundle, a smooth bundle (that is what ``the tangent bundle is a smooth bundle'' means: $\pi$ is smooth; \cref{def:tb-bundle}). A \textbf{vector field} on $M$ is a \emph{smooth section} of the tangent bundle: a smooth map $\sigma\colon M\to TM$ with
\[
  \pi\circ\sigma=\mathrm{id}_M
\]
(\cref{def:mf-section}). The set of all vector fields is denoted $\Gamma(TM)$, the standard notation for the set of all sections of a bundle (strictly one should write $\Gamma(TM,\pi,M)$, since all three data enter; this is short notation) \ts{11}{1}{2:11} \ts{11}{1}{2:59} \ts{11}{1}{3:19} \ts{11}{1}{3:40} \ts{11}{1}{6:06} \ts{11}{1}{6:27} \ts{11}{1}{7:08} \ts{11}{1}{7:29} \ts{11}{1}{8:18} \ts{11}{1}{8:59}.
\end{definition}
A vector field assigns to each point $p\in M$ an element $\sigma(p)\in T_pM$: because $\pi\circ\sigma=\mathrm{id}_M$ the vector provided at $p$ lies in the tangent space \emph{at $p$}, not in the tangent space at another point $q$; at another point another vector is provided. One could not simply write ``a map $M\to T_pM$'': one does not know at which
point one is, and one would have to say what it means for a vector field to be smooth, that is how one differentiates vector fields, which would need a new notion of differentiation. Instead, $M$ and $TM$ are smooth manifolds and one knows what a smooth map between smooth manifolds is
(\cref{ch:l07}): this is the whole idea of the definition \ts{11}{1}{4:22} \ts{11}{1}{4:43} \ts{11}{1}{5:04} \ts{11}{1}{5:24} \ts{11}{1}{5:45}.
From now on ``vector field'' always means smooth vector field; on a topological manifold one could only speak of continuous ones (a continuous map into the then topological manifold $TM$), and at whatever level of structure one is, the corresponding notion of field is obtained in the same way \ts{11}{1}{6:27} \ts{11}{1}{6:47}.

\begin{coursenote}[Notation $\Gamma$ asked in class]
A student asked whether $\Gamma(TM)$ is ``the collection of all $\sigma$'': yes, the set of all sections, those maps with $\pi\circ\sigma=\mathrm{id}_M$; the symbol carries the total space only because the bundle map and the base are understood. That the collection is a set at all (not a class) is clear since
sections are defined as maps, which are subsets of Cartesian products. \ts{11}{1}{8:18}
\end{coursenote}

\paragraph{Two operations on $\Gamma(TM)$.} The first one is harmless. For the time being call it $\oplus$:
\[
  \oplus\colon\Gamma(TM)\times\Gamma(TM)\to\Gamma(TM),\qquad(\sigma\oplus\tau)(p):=\sigma(p)+\tau(p),
\]
defined \textbf{pointwise}: at each $p\in M$ one adds the two vectors in the vector space $T_pM$ (the $+$ on the right is the one of $T_pM$, which makes sense because of the section property). No secret here: take two vector fields, look at every point and add the two vectors there; the collection of results is again a vector field
\ts{11}{1}{9:20} \ts{11}{1}{9:41} \ts{11}{1}{10:01} \ts{11}{1}{10:21}.

The second one is more interesting. It is a kind of scalar multiplication, but one may ask by what. One can of course multiply a vector field by a real number, but ``let us not be so childish'': one can multiply it by an entire \emph{function} on the manifold,
\[
  \odot\colon C^\infty(M)\times\Gamma(TM)\to\Gamma(TM),\qquad(f\odot\sigma)(p):=f(p)\cdot\sigma(p),
\]
again pointwise, with the scalar multiplication of $T_pM$ on the right. The function has different values at different points and scales the vector of the field at each point by the value it has \emph{there}; there is no need for one global real number that doubles or triples the whole field. ``It is very innocent, and it has
dramatic effects in what follows'' \ts{11}{1}{10:42} \ts{11}{1}{11:03} \ts{11}{1}{11:24} \ts{11}{1}{11:45} \ts{11}{1}{12:05} \ts{11}{1}{12:26}.
\begin{coursenote}[Smoothness of the results]
That $\sigma\oplus\tau$ and $f\odot\sigma$ are again \emph{smooth} sections is not discussed in class; in the induced charts of $TM$ (\cref{def:tb-chart}) the components of $\sigma\oplus\tau$ and $f\odot\sigma$ are $\sigma^a+\tau^a$ and $f\sigma^a$, smooth if $f,\sigma^a,\tau^a$ are. Not said in class.
\end{coursenote}

The set $C^\infty(M)$ has appeared in several forms (\cref{def:tv-cinf}). As a real vector space, multiplied by real numbers (the small dot). But here something else is meant: the multiplication of two functions, pointwise, which (\cref{ex:tv-cinf-alg}) makes $C^\infty(M)$ a \emph{ring}; two functions can also be added pointwise. So the structure on $\Gamma(TM)$ is an
addition and a multiplication, but the scalars, the elements of $C^\infty(M)$, are \emph{not the elements of a field, only of a ring}. That could seem harmless, and it is not \ts{11}{1}{12:26} \ts{11}{1}{12:46} \ts{11}{1}{13:07} \ts{11}{1}{13:28} \ts{11}{1}{13:48} \ts{11}{1}{14:29} \ts{11}{1}{14:51}.

% ==================================================================
\section{Rings and modules}\label{sec:rm-rings}

\begin{recall}[Rings]
The definition of ring and the examples were given in \cref{def:te-ring,ex:te-rings}; here only the distinctions are needed again, in the form of the lecture: a ring has $+$ and $\cdot$ (both maps $R\times R\to R$), $(R,+)$ satisfies the commutative, associative, neutral and inverse laws, and the distributive law links the two.
\end{recall}
The \textbf{distributive law} was not written down last time (``a self-correcting mistake''): it links the two operations, $a\cdot(b+c)=a\cdot b+a\cdot c$ and the like. For a ring, in contrast to a field, the following may be missing, quite independently of each other (almost: see the third point) \ts{11}{1}{17:58} \ts{11}{1}{18:19} \ts{11}{1}{18:42}:
\begin{enumerate}
  \item the commutativity of $\cdot$: if it holds, the ring is a \textbf{commutative ring};
  \item the neutral element of $\cdot$: if it exists, the ring is a \textbf{unital} ring; only unital rings will be considered here, but there are rings without a unit;
  \item the multiplicative inverse of every non-zero element (``zero'' being the neutral element of $+$): if it exists, the ring is a \textbf{division ring}; a non-commutative structure which is otherwise a field is a division ring (a ``skew field''; the lecture recalls the term only vaguely: [?]).
\end{enumerate}
A division ring is necessarily unital, since the inverse is defined through the neutral element of the product: $a\cdot b=1$. A field is a commutative division ring \ts{11}{1}{19:02} \ts{11}{1}{19:23} \ts{11}{1}{19:43} \ts{11}{1}{20:04} \ts{11}{1}{20:25} \ts{11}{1}{20:45} \ts{11}{1}{21:06}.

\begin{proposition}[$C^\infty(M)$ is a commutative unital ring, not a division ring]\label{prop:rm-cinf-ring}
$\bigl(C^\infty(M),+,\bullet\bigr)$ with the pointwise operations is a commutative unital ring that is not a division ring \ts{11}{1}{21:26}.
\end{proposition}
\begin{proof}
\emph{Commutative}: the pointwise product inherits the commutativity of the product in $\R$.
\emph{Unital}: the neutral element is the constant function $1$ (every point goes to the number $1$), which is smooth, and $1\bullet f=f$.
\emph{Not a division ring}: let $f\ne0$ vanish at one point $p$ (without being the zero function, which is the neutral element of $+$). There is then no $g$ with $f\bullet g=1$, since $f(p)g(p)=0\ne1$:
the pointwise multiplication cannot ``get rid of the zero'' at $p$ \ts{11}{1}{21:47} \ts{11}{1}{22:08} \ts{11}{1}{22:29} \ts{11}{1}{22:49} \ts{11}{1}{23:09}.
\end{proof}
So one is immediately dealing with a commutative and unital ring, ``but not, and that is of key importance, a so-called division ring'' \ts{11}{1}{23:09} \ts{11}{1}{23:30}.
\begin{exercise}[Zero and non-zero functions]\label{ex:rm-zero}
Show that on $M=\R$ the function $f(x)=x$ has no inverse in $C^\infty(\R)$, although it is non-zero. Is $f$ invertible on $M=\R\setminus\{0\}$? Which functions in $C^\infty(M)$ \emph{are} invertible, for any $M$? \emph{Not said in class.}
\end{exercise}

\begin{definition}[New concepts: $R$-module]\label{def:rm-module}
Let $R$ be a ring. An \textbf{$R$-module} is a set $\mathcal M$ (in class written $M$, ``M like module'', ``I could choose a different letter'', although $M$ is the manifold; here $\mathcal M$) with an addition and a scalar multiplication
\[
  \oplus\colon\mathcal M\times\mathcal M\to\mathcal M,\qquad\odot\colon R\times\mathcal M\to\mathcal M,
\]
which satisfy the vector space axioms for $\oplus$ (commutative, associative, neutral element, inverse), the axioms for the interplay of the two (associativity $(r\cdot s)\odot m=r\odot(s\odot m)$ and the distributive laws), and where the scalars are taken from the ring $R$ instead of from a field \ts{11}{1}{23:52} \ts{11}{1}{24:14} \ts{11}{1}{24:35} \ts{11}{1}{24:55} \ts{11}{1}{25:16} \ts{11}{1}{25:37}.
\end{definition}
So it is a ``vector space over a ring'', which would be perfect terminology; but modules behave so differently from vector spaces that they received their own name \ts{11}{1}{25:37} \ts{11}{1}{25:58}.
\begin{coursenote}[Left and right modules, and the unit]
With the multiplication from the left one speaks of a \emph{left} $R$-module; one can define \emph{right} modules, multiplying from the right, and a bimodule if both are given compatibly: ``let us not go into this, it is not necessary'', and over a commutative ring such as $C^\infty(M)$ the distinction disappears. For a unital ring one also asks $1\odot m=m$ for all $m$; this axiom is used implicitly below
(``it is a unital ring, $v=1\cdot v$''). \ts{11}{1}{24:55} \ts{11}{1}{1:12:42}
\end{coursenote}
The motivating example: the set $\Gamma(TM)$ with the pointwise operations $\oplus$ and $\odot$ is a $C^\infty(M)$-module. ``The tricky thing in the module is not the module structure, it is the ring structure that underlies it'' \ts{11}{1}{26:19} \ts{11}{1}{26:40}.
Why not restrict to real numbers, which would give back a real vector space? One could, but then a basis would always be (at best) \emph{uncountable} even for simple examples, as shown below; roughly, one wants to provide the coefficients of a vector field as \emph{functions}, which change from point to point, as fields in the physicist's sense do \ts{11}{1}{16:36} \ts{11}{1}{16:57} \ts{11}{1}{17:17} \ts{11}{1}{17:37}.

% ==================================================================
\section{Modules need not have a basis}\label{sec:rm-basis}

The key fact: unlike a vector space, a module generically does \emph{not} have a basis \ts{11}{1}{27:00} \ts{11}{1}{27:23}. If the underlying ring is a division ring there is not much difference: then a module is, in all its structural properties, like a vector space (the lack of commutativity plays a role in some calculations, but not in the
structure) and this will be proven. But $C^\infty(M)$ is not a division ring, so $\Gamma(TM)$ in general has no basis, and there are geometric examples \ts{11}{1}{27:43} \ts{11}{1}{28:04} \ts{11}{1}{28:24} \ts{11}{1}{28:45} \ts{11}{1}{29:05} \ts{11}{1}{29:26}.

\begin{theorem}[Modules over a division ring have a basis]\label{thm:rm-basis}
Let $D$ be a division ring and $V$ a $D$-module. Then $V$ has a basis \ts{11}{1}{30:51} \ts{11}{1}{31:17} \ts{11}{1}{31:39}.
\end{theorem}
\begin{corollary}[Every vector space has a basis]\label{cor:rm-vs}
Every field is a division ring, hence every vector space has a basis (a Hamel basis, \cref{def:te-basis}) \ts{11}{1}{31:39} \ts{11}{1}{32:00}.
\end{corollary}
Here a \textbf{basis} of a module $V$ over a ring $R$ is, as for vector spaces, a subset $B\subseteq V$ such that every element of $V$ can be written \emph{uniquely} as a finite linear combination $\sum r^ie_i$ with $e_i\in B$, $r^i\in R$ \ts{11}{1}{37:39}. The theorem will be proven below; first two examples, one with and one without a basis.

\begin{example}[Vector fields on $\R^2$: a basis exists]\label{ex:rm-r2}
Let $M=\R^2$. The set $\Gamma(T\R^2)$ is a $C^\infty(\R^2)$-module; the ring is not a division ring, yet it \emph{has} a basis: ``I never said that without a division ring there is no basis, only that with one you are guaranteed to have one''. Namely there exist two vector \emph{fields} $E_1,E_2$ (not vectors), for example the constant fields
along the two axes, such that for every vector field $\sigma$ there are two coefficient functions with $\sigma=\sigma^1\odot E_1\oplus\sigma^2\odot E_2$: at each point one scales $E_1$ and $E_2$ differently to produce the value of $\sigma$ there. This is the standard intuition about vector fields and it is fine
\ts{11}{1}{32:41} \ts{11}{1}{33:03} \ts{11}{1}{33:26} \ts{11}{1}{33:47} \ts{11}{1}{34:07} \ts{11}{1}{34:28} \ts{11}{1}{34:49} \ts{11}{1}{35:11} \ts{11}{1}{35:32} \ts{11}{1}{35:54} \ts{11}{1}{36:15} \ts{11}{1}{36:36} \ts{11}{1}{36:56}.
\end{example}

\begin{example}[Vector fields on $S^2$: no basis]\label{ex:rm-s2}
Let $M=S^2$. A basis means that a decomposition with \emph{unique} coefficient functions exists (otherwise one only has a generating system). The claim is that $\Gamma(TS^2)$, as a $C^\infty(S^2)$-module, has \textbf{no basis}. It is due to a famous and very intuitive result: \emph{you cannot comb a hairy sphere}. Imagine the sphere with a hair at
every point, none of length zero (``I am an ape, fully haired''); try to comb it smoothly so that the hair gives a smooth vector field, nowhere zero: there is always at some point a curl (a parting), and there the field is not smooth, or has to vanish; with a second curl, the same. Formally: there is no everywhere non-zero and everywhere
smooth vector field on $S^2$ (the hairy ball theorem) \ts{11}{1}{36:56} \ts{11}{1}{37:16} \ts{11}{1}{37:39} \ts{11}{1}{37:59} \ts{11}{1}{38:20} \ts{11}{1}{38:41} \ts{11}{1}{39:02} \ts{11}{1}{39:22} \ts{11}{1}{39:43} \ts{11}{1}{40:03} \ts{11}{1}{40:24}.
How does this relate to a basis? If there is not even one everywhere non-zero smooth vector field, much less are there two that are linearly independent at every point, such that every vector field on the sphere could be uniquely decomposed in terms of them. A basis element cannot vanish; and if a candidate $E_1$ vanishes at some
point while $E_2$ happens to be non-zero there, a vector field with a component along a ``second direction'' at that point cannot be produced by $E_1$, since $E_1$ is zero there, and the coefficient function does not help. So the module over the non-division ring $C^\infty(S^2)$ has no basis: it is no surprise that here the module has no basis, whereas on $\R^2$, thanks to the special structure
of $\R^2$, everything was fine \ts{11}{1}{41:01} \ts{11}{1}{41:22} \ts{11}{1}{41:42} \ts{11}{1}{42:03} \ts{11}{1}{42:24} \ts{11}{1}{42:45} \ts{11}{1}{43:05}.
\end{example}
\begin{secondary}[A complete argument for the sphere]
\emph{Not said in class.} The lecture's argument is an intuition; here is a way to make it a proof, using the hairy ball theorem. Suppose $\Gamma(TS^2)$ had a basis $E_1,\dots,E_n$. Then every vector $v\in T_pS^2$ is a combination of the $E_i(p)$ (extend $v$ to a field with a bump function), so $E_i(p)$ span $T_pS^2$ for every $p$, and
$\Phi\colon S^2\times\R^n\to TS^2$, $(p,c)\mapsto\sum c^iE_i(p)$ is a surjective bundle map of constant rank. If $n>2$ its kernel is a vector bundle of positive rank, with a non-zero smooth section $c(p)$ (bump functions); then $\sum c^i\odot E_i=0$ with $c\ne0$ contradicts the uniqueness of the coefficients. So $n=2$,
$\Phi$ is an isomorphism, $TS^2$ is trivial and $E_1$ is an everywhere non-zero vector field, contradicting the hairy ball theorem.
\end{secondary}
\begin{coursenote}[Sets and sequences: bases]
At the end of the proof below the lecturer notes that a basis is really not a subset but more like a \emph{sequence} (a family): if the same element appears twice the elements are linearly dependent, whereas as a set $\{v,v\}=\{v\}$ is independent; and if one speaks of ``the first, second component'' an order is used, which a subset does not have. ``Don't get confused by this''; it
is a notational issue \ts{11}{1}{25:51} \ts{11}{1}{26:12} \ts{11}{1}{26:32}.
\end{coursenote}

\paragraph{Why not use real coefficients only.} One might retreat and consider $\Gamma(TM)$ simply as a real vector space, which it is; then it has a basis. But the coefficients in a decomposition would be real numbers. On $\R^2$, the pair $E_1,E_2$ would no longer be a basis: only fields with the \emph{same} combination of $E_1$ and $E_2$ at every point could be expressed, such as constant
fields. One would need uncountably many basis elements, which is not tractable. Therefore one wants coefficient \emph{functions} for tensor fields, and has to deal with this situation \ts{11}{1}{43:49} \ts{11}{1}{44:10} \ts{11}{1}{44:30} \ts{11}{1}{44:50}.
\medskip
Now the situation is understood algebraically; this amounts to the proof of \cref{thm:rm-basis}. The proof needs Zorn's lemma, the incarnation in which the \emph{axiom of choice} (\cref{sec:se-choice}) will be used. Everything that now comes is of real importance, the techniques and notions in it being ubiquitous in mathematics \ts{11}{1}{45:11} \ts{11}{1}{45:33} \ts{11}{1}{45:53} \ts{11}{1}{46:16}.

% ==================================================================
\section{Zorn's lemma}\label{sec:rm-zorn}

Zorn's lemma is in fact \emph{equivalent} to the axiom of choice in Zermelo--Fraenkel set theory (without choice): it is just another way to express it, but in a form that is very often useful. It is stated first, and the terms are explained afterwards \ts{11}{1}{46:16} \ts{11}{1}{46:37} \ts{11}{1}{46:57}.

\begin{theorem}[Zorn's lemma]\label{thm:rm-zorn}
Let $(P,\le)$ be a partially ordered set in which every totally ordered subset $T\subseteq P$ has an upper bound in $P$. Then $P$ contains a maximal element \ts{11}{1}{47:18} \ts{11}{1}{47:41} \ts{11}{1}{48:01} \ts{11}{1}{48:22}.
\end{theorem}
The proof is not given; in ZF set theory (without choice) the axiom of choice is equivalent to Zorn's lemma \ts{11}{1}{59:36} \ts{11}{1}{59:58}.
(``Apart from four undefined terms it is very enlightening'': partially ordered set, totally ordered subset, upper bound, maximal element.)

\begin{definition}[New concepts: partially ordered set]\label{def:rm-poset}
A \textbf{partially ordered set} is a set $P$ with a relation $\le$ (``keep an open mind about what this less-or-equal is'': it is not necessarily about numbers) which is
\begin{enumerate}
  \item \emph{reflexive}: $a\le a$ for all $a\in P$;
  \item \emph{antisymmetric} (not symmetric, as for equivalence relations): $a\le b$ and $b\le a$ imply $a=b$;
  \item \emph{transitive}: $a\le b$ and $b\le c$ imply $a\le c$ \ts{11}{1}{48:43} \ts{11}{1}{49:05} \ts{11}{1}{49:26} \ts{11}{1}{49:47} \ts{11}{1}{50:07}.
\end{enumerate}
\end{definition}
\begin{definition}[New concepts: totally ordered subset]\label{def:rm-total}
A subset $T\subseteq P$ is \textbf{totally ordered} if the restriction of $\le$ to $T$ satisfies: (a) \emph{totality}: for all $a,b\in T$, $a\le b$ or $b\le a$ (inclusive ``or''); (b) antisymmetry; (c) transitivity. Reflexivity is not required separately: it follows from totality
(take $a=b$) \ts{11}{1}{50:28} \ts{11}{1}{50:51} \ts{11}{1}{51:11} \ts{11}{1}{51:35} \ts{11}{1}{51:56} \ts{11}{1}{52:17}.
\end{definition}
In a partially ordered set two elements can stand in no relation to each other; in a totally ordered one any two elements are related one way or the other \ts{11}{1}{52:37} \ts{11}{1}{52:58}. A student asks whether the complex numbers are partially ordered. A partial order on $\C$ can certainly be defined: relate only objects of the real line by the standard $\le$ (nothing else is related). That is reflexive,
antisymmetric and transitive, hence a partial order, though not a total one on $\C$; restricted to $\R$ it is a total order since any two real numbers are related \ts{11}{1}{53:25} \ts{11}{1}{53:47} \ts{11}{1}{54:07} \ts{11}{1}{54:28} \ts{11}{1}{54:49}.

\begin{definition}[New concepts: upper bound, maximal element]\label{def:rm-upper}
Let $T\subseteq P$. An element $u\in P$ is an \textbf{upper bound} of $T$ if $t\le u$ for all $t\in T$ (the relation is that of $P$; $u$ comes from $P$, not necessarily from $T$). An element $m\in P$ is \textbf{maximal} if there is no $x\in P$ with $m\le x$
\ts{11}{1}{55:10} \ts{11}{1}{55:31} \ts{11}{1}{55:52} \ts{11}{1}{56:16} \ts{11}{1}{56:42} \ts{11}{1}{57:06} \ts{11}{1}{57:31}.
\end{definition}
\begin{coursenote}[Maximal element: a slip]
As written, with $x$ arbitrary, the definition of maximal is wrong: reflexivity gives $m\le m$, so no element would be maximal; the right statement is ``there is no $x\in P$ with $m\le x$ and $x\ne m$'' (equivalently: $m\le x$ implies $x=m$). Not said in class. The remark that follows in class is correct: the tempting alternative ``$u\le m$ for all $u\in P$'' is a different notion
(a \emph{greatest} element) and in a partial order is stronger: it moves the negation incorrectly, since a negation in front of $\le$ can be flipped only when any two elements are related, which holds for a total, not for a partial order. In the class the lecturer found the same error in his own preparation
(``very honest''): with a partial order there can be elements of $P$ unrelated to $m$ \ts{11}{1}{57:52} \ts{11}{1}{58:12} \ts{11}{1}{58:33} \ts{11}{1}{58:54} \ts{11}{1}{59:15}.
\end{coursenote}

\paragraph{The student question on upper bounds.} A student asks whether the upper bound $u$ has to be in $T$ or may be anywhere in $P$ (and raises the case that $u$ may then be unrelated to the elements of $T$). The lecturer first tries the modified definition ($u\in T$, to play it safe), checks whether it suits the application below (where the upper bound is the union of the chain, and he is not sure it lies in $T$ [?]), draws a diagram and thereupon \emph{returns to the original definition}. The reasoning: in a poset $P$ with a chain $T$ and some element $u\in P$ unrelated to $T$, then $u$ is simply not an upper bound of $T$ (the condition $t\le u$ for all $t$ fails), but
another element which is related to every $t\in T$ is, and that is no problem; the requirement ``$T$ totally ordered'' ensures that elements of $T$ can be compared with the bound. The original definition (with $u\in P$) is the \emph{weaker} hypothesis (more chains have upper bounds), hence Zorn's lemma is stronger and in the application one wants the freedom to
have the upper bound outside $T$ \ts{11}{1}{1:00:22} \ts{11}{1}{1:00:45} \ts{11}{1}{1:01:06} \ts{11}{1}{1:01:33} \ts{11}{1}{1:01:59} \ts{11}{1}{1:02:19} \ts{11}{1}{1:02:40} \ts{11}{1}{1:03:01} \ts{11}{1}{1:03:22} \ts{11}{1}{1:03:43} \ts{11}{1}{1:04:07} \ts{11}{1}{1:04:30} \ts{11}{1}{1:04:51} \ts{11}{1}{1:05:15} \ts{11}{1}{1:05:36}
\ts{11}{1}{1:05:57} \ts{11}{1}{1:06:18} \ts{11}{1}{1:06:38} \ts{11}{1}{1:06:59} \ts{11}{1}{1:07:20} \ts{11}{1}{1:07:44} \ts{11}{1}{1:08:04} \ts{11}{1}{1:08:25} \ts{11}{1}{1:08:46} \ts{11}{1}{1:09:11} \ts{11}{1}{1:09:31} \ts{11}{1}{1:09:53}.

% ==================================================================
\section{Proof that modules over a division ring have a basis}\label{sec:rm-proof}

\begin{proof}[Proof of \cref{thm:rm-basis}]
\emph{Step 1: a generating system.} Let $S\subseteq V$ be a \textbf{generating system}: for every $v\in V$ there are finitely many $e_1,\dots,e_n\in S$ and $v^1,\dots,v^n\in D$ with
\[
  v=v^ie_i.
\]
This is no basis yet: nothing is said about uniqueness. A generating system exists: take $S=V$, since $v=1\cdot v$ in the unital ring $D$. It may be ``fantastically huge'', but the question is whether a basis can be extracted from it \ts{11}{1}{1:10:16} \ts{11}{1}{1:10:37} \ts{11}{1}{1:10:59} \ts{11}{1}{1:11:20} \ts{11}{1}{1:11:40} \ts{11}{1}{1:12:00} \ts{11}{1}{1:12:20} \ts{11}{1}{1:12:42} \ts{11}{1}{1:13:02} \ts{11}{1}{1:13:24}.

\emph{Step 2: a partially ordered set.} Let
\[
  P:=\{U\in\mathcal P(S)\mid U\text{ linearly independent}\},
\]
partially ordered by inclusion $\subseteq$ (two subsets either are contained in one another or neither is; no problem, only a partial order is needed). Linear independence means that every finite subset $\{e_1,\dots,e_n\}$ of $U$ satisfies $a^ie_i=0\Rightarrow a^1=\dots=a^n=0$ \ts{11}{1}{1:13:45} \ts{11}{1}{1:14:08} \ts{11}{1}{1:14:29} \ts{11}{1}{1:14:50} \ts{11}{1}{1:15:12} \ts{11}{1}{1:15:33} \ts{11}{1}{1:15:53}.

\emph{Step 3: Zorn.} Let $T\subseteq P$ be an \emph{arbitrary} totally ordered subset (no further specification, to be able to apply Zorn's lemma). Then the union $\bigcup T$ of its elements is an upper bound of $T$ (every member of $T$ is contained in it) and is linearly independent: a finite subset of $\bigcup T$ lies in one member of $T$ (the members are totally
ordered by inclusion, so the largest of the finitely many members containing the elements is one of them), which is independent. So $\bigcup T\in P$. By Zorn's lemma $P$ has a maximal element $B$ (possibly several; one is picked: here the axiom of choice enters). By construction $B$ is a \emph{maximal linearly independent subset} of $S$ in the sense of inclusion: nothing can be
added without losing independence \ts{11}{1}{1:16:14} \ts{11}{1}{1:16:36} \ts{11}{1}{1:16:57} \ts{11}{1}{1:17:17} \ts{11}{1}{1:17:38} \ts{11}{1}{1:17:58} \ts{11}{1}{1:18:19} \ts{11}{1}{1:18:40} \ts{11}{1}{1:19:02} \ts{11}{1}{1:19:23} \ts{11}{1}{1:19:43} \ts{11}{1}{1:20:05} \ts{11}{1}{1:20:26} \ts{11}{1}{1:20:47}.

\emph{Step 4: $B$ spans $S$.} The claim is $S\subseteq\mathrm{span}_DB$ (finite linear combinations with coefficients in $D$) \ts{11}{1}{1:21:08} \ts{11}{1}{1:21:30} \ts{11}{1}{1:21:50}. Let $v\in S$ be arbitrary. If $v\in B$ there is nothing to show. If not, $B\cup\{v\}$ strictly contains the maximal $B$, so it is linearly dependent: there are
$e_1,\dots,e_n\in B$ and $a^1,\dots,a^n,b\in D$, not all zero, with
\[
  a^ie_i+b\,v=0.
\]
Here $b\ne0$: if $b=0$ then $a^ie_i=0$ with not all $a^i$ zero, contradicting the linear independence of $B$. Since $D$ is a division ring, $b$ has a multiplicative inverse, and
\[
  v=-\bigl(b^{-1}a^i\bigr)e_i,
\]
a finite linear combination of elements of $B$ (``this is where we use the division ring''; for non-commutative $D$ the order of the factors matters and the scalars act from the left) \ts{11}{1}{1:22:14} \ts{11}{1}{1:22:36} \ts{11}{1}{1:22:57} \ts{11}{1}{1:23:17} \ts{11}{1}{1:23:41} \ts{11}{1}{1:24:03} \ts{11}{1}{1:24:26} \ts{11}{1}{1:24:46} \ts{11}{1}{1:25:07} \ts{11}{1}{1:25:31} \ts{11}{1}{1:27:12} \ts{11}{1}{1:27:33} \ts{11}{1}{1:27:54} \ts{11}{1}{1:28:16} \ts{11}{1}{1:28:36} \ts{11}{1}{1:28:58} \ts{11}{1}{1:29:20} \ts{11}{1}{1:29:45} \ts{11}{1}{1:30:06} \ts{11}{1}{1:30:26} \ts{11}{1}{1:30:49}.

\emph{Step 5: conclusion.} $V$ is spanned by its generating system $S$, and $S$ is spanned by $B$, so $B$ generates $V$; it is linearly independent, hence unique decompositions exist (the difference of two decompositions is a vanishing combination). Thus $B$ is a basis of $V$, ``a basis, not the unique basis''; Zorn's lemma only guarantees that some maximal element exists
\ts{11}{1}{1:31:09} \ts{11}{1}{1:31:31} \ts{11}{1}{1:31:51}.
\end{proof}
\begin{coursenote}[Details left out in class]
(i) The empty chain $T=\emptyset$ has $\emptyset\in P$ as an upper bound: this is why Zorn's lemma in the form above needs $P\ne\emptyset$, which is guaranteed since $\emptyset\in P$. (ii) The lecturer's ``union of independent sets'' is exactly the finite-subset argument given in Step 3. (iii) The case $v\in B$ in Step 4 was discussed in class in
the form that $B\cup\{v\}=B$ is still independent, whence the remark on sets and sequences above. Not said in class. \ts{11}{1}{1:25:31}
\end{coursenote}
So every vector space has a Hamel basis (the construction may use infinite sets, but one speaks of finite linear combinations; again for vector spaces \emph{without} further structure such as a topology or a norm), and more generally every module over a division ring has a basis. But $C^\infty(M)$ is not a division ring, hence $\Gamma(TM)$ regarded as $C^\infty(M)$-module is
\emph{not guaranteed} to have a basis: it has one for $\R^2$, none for $S^2$ \ts{11}{1}{1:32:13} \ts{11}{1}{1:32:34} \ts{11}{1}{1:32:58} \ts{11}{1}{1:33:22} \ts{11}{1}{1:33:44}.

\begin{intuition}[No global component functions]\label{int:rm-nocomp}
If one thinks of vectors and vector fields always in terms of their components, one has the illusion that there are always components; but components refer to a basis, and if the module of vector fields has no basis there are \emph{no component functions}. One must stop thinking of tensor fields, in their most general form, in terms of component functions: ``there
are no such component functions in general, full stop''. Locally this is possible: on a patch of the sphere there are component functions, one changes chart and gets component functions on every chart, and that is what is done all the time; but there is not necessarily a global one, and this has nothing to do with choosing charts: abstractly no such thing exists. The formalism, which looked like formalism for
formalism's sake, has a geometric meaning \ts{11}{1}{1:34:04} \ts{11}{1}{1:34:25} \ts{11}{1}{1:34:45} \ts{11}{1}{1:35:07}.
\end{intuition}

% ==================================================================
\section{Free and projective modules}\label{sec:rm-projective}

Some terminology, which emerges here \ts{11}{1}{1:35:07}.

\begin{definition}[New concepts: free module]\label{def:rm-free}
An $R$-module is called \textbf{free} if it has a basis. This does not require $R$ to be a division ring: for example $\Gamma(T\R^2)$ is free, $\Gamma(TS^2)$ is not free. Unless $R$ is a division ring a module over it is not guaranteed to be free; sometimes it is, sometimes it is not \ts{11}{1}{1:35:28} \ts{11}{1}{1:35:50} \ts{11}{1}{1:36:11}.
\end{definition}
Free modules can be described via exact and split exact sequences, which will not be gone into here (``we may come back to it later in terms of cohomology''): a free module has a certain property in terms of such sequences, but there is another type of module that shares exactly this property without being free \ts{11}{1}{1:36:11} \ts{11}{1}{1:36:32} \ts{11}{1}{1:36:53}.
\begin{definition}[New concepts: projective module]\label{def:rm-proj}
An $R$-module $\Gamma$ is called \textbf{projective} if it is a direct summand of a free module $F$: there is an $R$-module $Q$ such that
\[
  \Gamma\oplus Q\ \text{ is free}\ (=F).
\]
(The name is not explained in the lecture; it comes from the exact-sequence property just mentioned, not said in class.) Clearly every free module is projective (take $Q=\{0\}$): \emph{free implies projective} \ts{11}{1}{1:36:53} \ts{11}{1}{1:37:15} \ts{11}{1}{1:37:37} \ts{11}{1}{1:37:58} \ts{11}{1}{1:38:18}.
\end{definition}
Projective modules generalise free modules. A finitely generated free $R$-module $F$ is isomorphic to a direct sum of copies of the ring, $F\cong R\oplus\dots\oplus R$, exactly as for vector spaces: a finite dimensional vector space is the direct sum of $\dim$ many copies of the field. (As a set the direct sum is a Cartesian product, with the linear structure working
componentwise, as known from the problem sheet.) A projective module fails to be free when the summand $Q$ cannot be chosen trivial: only after adding another module is it free \ts{11}{1}{1:38:38} \ts{11}{1}{1:39:01} \ts{11}{1}{1:39:21} \ts{11}{1}{1:39:42} \ts{11}{1}{1:40:03}.

The reason for all of this is a theorem, quoted without proof; the lecture attributes it to Swan and others, with extensions by others (the \emph{Serre--Swan theorem}).
\begin{theorem}[Serre--Swan, in the form quoted]\label{thm:rm-serre-swan}
The set of all smooth sections of a vector bundle $E\to M$ over a smooth manifold $M$ (a bundle whose fibres are vector spaces) is a \emph{finitely generated projective} $C^\infty(M)$-module, denoted $\Gamma(E)$ \ts{11}{1}{1:40:27} \ts{11}{1}{1:40:47} \ts{11}{1}{1:41:11} \ts{11}{1}{1:41:31} \ts{11}{1}{1:41:52} \ts{11}{1}{1:42:19}.
\end{theorem}
Examples of vector bundles: the tangent bundle, the cotangent bundle (every fibre the cotangent space), a tensor bundle (at every point a tensor space, itself a vector space): ``it is very general''. So $\Gamma(E)\oplus Q$ is free for some $C^\infty(M)$-module $Q$; the module $\Gamma(E)$ may have no basis, but behaves as a summand of a module that very much behaves like a vector space.
If $Q$ could be chosen $\{0\}$ (the module with only the additive zero) it would be free: for $E=T\R^2$ it can, for $E=TS^2$ it cannot \ts{11}{1}{1:42:41} \ts{11}{1}{1:43:01} \ts{11}{1}{1:43:22}.
\begin{coursenote}[Hypotheses of the theorem]
As said in class (not checked): the theorem is attributed to Swan. R.~G.~Swan, \emph{Vector bundles and projective modules}, Trans.\ Amer.\ Math.\ Soc.\ 105 (1962), proved it for vector bundles over a compact Hausdorff space and continuous sections; the version for smooth manifolds with $C^\infty(M)$ is the one quoted here. The
finite generation uses that a vector bundle over a manifold is covered by finitely many trivialising open sets (``of finite type''), true for manifolds; the precise hypotheses are not discussed in the lecture [?]. Not said in class.
\end{coursenote}
So $Q$ quantifies how much $\Gamma(E)$ fails to have a basis (``corollary, if you want''). Since a module over a non-division ring is not expected to have a basis, when one has one it is an accident, a coincidence; whoever is surprised that the vector fields on the sphere have no basis can say: ``why would it? It is just a projective module, not a free one'': the algebraic expression of the geometric fact,
which is of course not special to the sphere \ts{11}{1}{1:43:22} \ts{11}{1}{1:43:43} \ts{11}{1}{1:44:06} \ts{11}{1}{1:44:28} \ts{11}{1}{1:44:48}.

\paragraph{Homomorphisms of projective modules.} Two more remarks to wrap up \ts{11}{1}{1:44:48}.
\begin{theorem}[Hom of projective modules]\label{thm:rm-hom}
Let $P,Q$ be projective $R$-modules over a \emph{commutative} ring $R$. Then $\mathrm{Hom}_R(P,Q)$, the set of all $R$-linear maps $P\to Q$ (linear as for vector spaces, with respect to the linear structure of the modules and the ring), is again projective; if $P$ and $Q$ are finitely generated, so is $\mathrm{Hom}_R(P,Q)$ \ts{11}{1}{1:45:09} \ts{11}{1}{1:45:29} \ts{11}{1}{1:45:50} \ts{11}{1}{1:46:11} \ts{11}{1}{1:46:31}.
\end{theorem}
The proof is not given; the only surprise is the commutativity. Equip $\mathrm{Hom}_R(P,Q)$ with the pointwise addition and the multiplication by ring elements, $(r\odot\varphi)(v):=r\odot\varphi(v)$. For this to be again $R$-linear one needs $r\odot\varphi(s\odot v)=s\odot(r\odot\varphi(v))$: by linearity of $\varphi$ the
$s$ comes out, $\varphi(s\odot v)=s\odot\varphi(v)$, but then to bring $r$ and $s$ in the right order one needs $rs=sr$: ``only over a commutative ring does one have this actually expected structure'', homomorphisms are again a structure of the same type, and projectivity and finite generation are preserved
\ts{11}{1}{1:46:52} \ts{11}{1}{1:47:12} \ts{11}{1}{1:47:36} \ts{11}{1}{1:47:57} \ts{11}{1}{1:48:17} \ts{11}{1}{1:48:38} \ts{11}{1}{1:48:59} \ts{11}{1}{1:49:20}.
\begin{coursenote}[Dots in the lecture's computation]
The board computation of this commutativity argument contains several dots ($\odot$ in the ring module, the ring product, the scalar multiplication in $Q$) and the lecturer himself says ``oh dear, we have various dots'' and ``I wrote the same thing twice, not very meaningful''; the above is the reconstructed content. Not said in class. \ts{11}{1}{1:47:57}
\end{coursenote}

In particular one can apply the theorem with $R=C^\infty(M)$, which is commutative, and send vector fields to $C^\infty(M)$ itself (trivially a module over itself). The result
\[
  \Gamma(TM)^*:=\mathrm{Hom}_{C^\infty(M)}\bigl(\Gamma(TM),C^\infty(M)\bigr)
\]
is again a finitely generated module over $C^\infty(M)$, the \textbf{dual} of the module of vector fields, and one can show that it is the module of sections of the cotangent bundle, $\Gamma(TM)^*=\Gamma(T^*M)$ (not proven). That closes the loop \ts{11}{1}{1:49:20} \ts{11}{1}{1:49:41} \ts{11}{1}{1:50:02} \ts{11}{1}{1:50:24}.

% ==================================================================
\section{Tensor fields}\label{sec:rm-tensorfields}

One is now fully prepared for the standard textbook definition \ts{11}{1}{1:50:24} \ts{11}{1}{1:50:45}.
\begin{definition}[New concepts: tensor field]\label{def:rm-tensorfield}
A \textbf{$(r,s)$ tensor field} $T$ on a smooth manifold $M$ (in the lecture $(r,s)$; $(p,q)$ in \cref{def:te-tensor}) is a $C^\infty(M)$-\emph{multilinear} map
\[
  T\colon\underbrace{\Gamma(T^*M)\times\dots\times\Gamma(T^*M)}_{r\text{ copies}}\times\underbrace{\Gamma(TM)\times\dots\times\Gamma(TM)}_{s\text{ copies}}\longrightarrow C^\infty(M).
\]
It eats $r$ covector fields and $s$ vector fields and makes of them a smooth function on the manifold \ts{11}{1}{1:51:07} \ts{11}{1}{1:51:31} \ts{11}{1}{1:51:52} \ts{11}{1}{1:52:12}.
\end{definition}
Pointwise: at a given point the tensor field takes $r$ covectors, those given by the covector fields, and $s$ vectors, and returns a real number; doing this at every point gives a function, and that everything is smooth, and that multilinearity holds, is encoded in the fact that one has a $C^\infty(M)$-multilinear map. The set of all tensor fields is again
a module over $C^\infty(M)$ \ts{11}{1}{1:52:12} \ts{11}{1}{1:52:34} \ts{11}{1}{1:52:54} \ts{11}{1}{1:53:16}.
\begin{remark}[$C^\infty(M)$-multilinear, not $\R$-multilinear]\label{rem:rm-multilinear}
Textbooks write ``multilinear'' and sometimes say it is not $\R$-multilinear but $C^\infty(M)$-multilinear, and \emph{that is the key thing to check}: since the objects one feeds in are elements of a module, multilinearity must refer to the ring structure, otherwise the definition makes no sense
(compare \cref{ex:rm-leibniz}). One now sees the module structure ``lurking underneath'' the textbook definition \ts{11}{1}{1:53:16} \ts{11}{1}{1:53:37} \ts{11}{1}{1:53:57} \ts{11}{1}{1:54:18}.
\end{remark}
\begin{exercise}[Linear over $\R$ but not over $C^\infty(M)$]\label{ex:rm-leibniz}
Show that $\Gamma(TM)\to C^\infty(M)$, $X\mapsto X(g)$ for a fixed $g\in C^\infty(M)$ is $C^\infty(M)$-linear, whereas $f\mapsto X(f)$ (with fixed $X$, defined on the \emph{ring} $C^\infty(M)$) is only $\R$-linear, not $C^\infty(M)$-linear: the Leibniz rule (\cref{def:tv-derivation}) prevents it.
\emph{Not said in class.}
\end{exercise}
Hence one would never get the idea that the set of tensor fields should itself have a basis: not even the objects that one feeds in have necessarily a basis, since everything is a module structure and not a vector space structure. Keeping that in mind one can of course work \emph{locally} with a local basis, but not with a global one. This is the relationship of modules,
vector fields and tensor fields \ts{11}{1}{1:54:18} \ts{11}{1}{1:54:39} \ts{11}{1}{1:54:59}.

% ==================================================================
\section*{Key concepts}
\begin{keyconcepts}
\textbf{Retrieval questions.}
\begin{enumerate}
  \item What is a vector field, and how do the pointwise $\oplus$ and the multiplication by a function work on $\Gamma(TM)$? Why is a function (not only a real number) the right scalar?
  \item Which of the ring properties (commutative, unital, division) does $C^\infty(M)$ have? Why is it not a division ring?
  \item Define an $R$-module. What is the difference from a vector space, and what is the key fact about bases?
  \item State Zorn's lemma, with partial order, totally ordered subset, upper bound and maximal element. Which formulation of ``maximal'' is correct for a partial order?
  \item Prove that a module over a division ring has a basis: which poset, which chain, why is the union an upper bound, where is the division ring used?
  \item Why does $\Gamma(T\R^2)$ have a basis and $\Gamma(TS^2)$ not? Why can one not think of tensor fields globally in terms of components?
  \item What are free and projective modules? What does the theorem on sections of vector bundles say, and what does $Q$ measure?
  \item Why must a tensor field be $C^\infty(M)$-multilinear? What is $\Gamma(TM)^*$?
\end{enumerate}
\textbf{Concepts.}
\begin{itemize}
  \item Vector field, $\Gamma(TM)$, pointwise operations (\cref{def:rm-field}); ring $C^\infty(M)$ (\cref{prop:rm-cinf-ring}); module (\cref{def:rm-module}).
  \item Basis of a module over a division ring, Zorn's lemma (\cref{thm:rm-basis,cor:rm-vs,thm:rm-zorn,def:rm-poset,def:rm-total,def:rm-upper}); $\R^2$ versus $S^2$ (\cref{ex:rm-r2,ex:rm-s2,int:rm-nocomp}).
  \item Free and projective modules, sections of vector bundles, $\mathrm{Hom}$ of projective modules (\cref{def:rm-free,def:rm-proj,thm:rm-serre-swan,thm:rm-hom}).
  \item Tensor fields as $C^\infty(M)$-multilinear maps (\cref{def:rm-tensorfield,rem:rm-multilinear}).
\end{itemize}
\end{keyconcepts}

\flushside
\end{paracol}
